\documentclass[10pt,a4paper]{amsart}
\usepackage{amsmath,amssymb,amsthm,graphicx}
% Retained minimal preamble for the Stacks Project chapter more-algebra.tex (More on Algebra).
% The chapter file pinned-original-source.tex carries only \input{preamble} on its first line, so no
% preamble text lives inside the chapter itself.  The Stacks Project uses one single shared preamble for
% every chapter; the snapshot stages it identically as sources/stacks-chow/preamble.tex and
% sources/stacks-etale-cohomology/preamble.tex.  The block below is copied mechanically, line for line,
% from sources/stacks-chow/preamble.tex lines 14-18 (the Xy-pic setup) and 154-217 (the native theorem
% declarations together with the mathematical macros).  Nothing was added, removed or redefined; the
% wrapper class and the amsmath/amssymb/amsthm/graphicx base packages are supplied by the assembler.
% Chapter: more-algebra.tex -- More on Algebra.
% For commutative diagrams we use Xy-pic
\usepackage[all]{xy}

% We use 2cell for 2-commutative diagrams.
\xyoption{2cell}

% Theorem environments.
%
\theoremstyle{plain}
\newtheorem{theorem}[subsection]{Theorem}
\newtheorem{proposition}[subsection]{Proposition}
\newtheorem{lemma}[subsection]{Lemma}

\theoremstyle{definition}
\newtheorem{definition}[subsection]{Definition}
\newtheorem{example}[subsection]{Example}
\newtheorem{exercise}[subsection]{Exercise}
\newtheorem{situation}[subsection]{Situation}

\theoremstyle{remark}
\newtheorem{remark}[subsection]{Remark}
\newtheorem{remarks}[subsection]{Remarks}

\numberwithin{equation}{subsection}

% Macros
%
\def\lim{\mathop{\mathrm{lim}}\nolimits}
\def\colim{\mathop{\mathrm{colim}}\nolimits}
\def\Spec{\mathop{\mathrm{Spec}}}
\def\Hom{\mathop{\mathrm{Hom}}\nolimits}
\def\Ext{\mathop{\mathrm{Ext}}\nolimits}
\def\SheafHom{\mathop{\mathcal{H}\!\mathit{om}}\nolimits}
\def\SheafExt{\mathop{\mathcal{E}\!\mathit{xt}}\nolimits}
\def\Sch{\mathit{Sch}}
\def\Mor{\mathop{\mathrm{Mor}}\nolimits}
\def\Ob{\mathop{\mathrm{Ob}}\nolimits}
\def\Sh{\mathop{\mathit{Sh}}\nolimits}
\def\NL{\mathop{N\!L}\nolimits}
\def\CH{\mathop{\mathrm{CH}}\nolimits}
\def\proetale{{pro\text{-}\acute{e}tale}}
\def\etale{{\acute{e}tale}}
\def\QCoh{\mathit{QCoh}}
\def\Ker{\mathop{\mathrm{Ker}}}
\def\Im{\mathop{\mathrm{Im}}}
\def\Coker{\mathop{\mathrm{Coker}}}
\def\Coim{\mathop{\mathrm{Coim}}}

% Boxtimes
%
\DeclareMathSymbol{\boxtimes}{\mathbin}{AMSa}{"02}

%
% Macros for moduli stacks/spaces
%
\def\QCohstack{\mathcal{QC}\!\mathit{oh}}
\def\Cohstack{\mathcal{C}\!\mathit{oh}}
\def\Spacesstack{\mathcal{S}\!\mathit{paces}}
\def\Quotfunctor{\mathrm{Quot}}
\def\Hilbfunctor{\mathrm{Hilb}}
\def\Curvesstack{\mathcal{C}\!\mathit{urves}}
\def\Polarizedstack{\mathcal{P}\!\mathit{olarized}}
\def\Complexesstack{\mathcal{C}\!\mathit{omplexes}}
% \Pic is the operator that assigns to X its picard group, usage \Pic(X)
% \Picardstack_{X/B} denotes the Picard stack of X over B
% \Picardfunctor_{X/B} denotes the Picard functor of X over B
\def\Pic{\mathop{\mathrm{Pic}}\nolimits}
\def\Picardstack{\mathcal{P}\!\mathit{ic}}
\def\Picardfunctor{\mathrm{Pic}}
\def\Deformationcategory{\mathcal{D}\!\mathit{ef}}

% --- appended 2026-10-12 (Chief): the project own definitions of the reference/slogan/history
% display environments, copied mechanically from the shared Stacks preamble lines 8-12
% (\usepackage{verbatim} plus the three \newenvironment lines, which the chapter preamble keeps
% verbatim from the project preamble). Nothing added, removed or redefined: these environments
% carry no mathematical content and the project itself typesets them as comments. Without these
% three lines a counted statement whose own text uses them fails with Environment slogan undefined.
\usepackage{verbatim}
\newenvironment{reference}{\comment}{\endcomment}
\newenvironment{slogan}{\comment}{\endcomment}
\newenvironment{history}{\comment}{\endcomment}
\usepackage{cleveref}
\crefname{originallemma}{original lemma}{original lemmas}
\Crefname{originallemma}{Original lemma}{Original lemmas}
\makeatletter
\newlabel{lemma-geometric-branches}{{0C25}{source}{}{sourcepointer-1}{}}
\newlabel{lemma-geometric-branches@cref}{{[originallemma][0][]0C25}{[1][0][]source}{}{}{}}
\makeatother
\crefname{originallemma}{original lemma}{original lemmas}
\Crefname{originallemma}{Original lemma}{Original lemmas}
\makeatletter
\newlabel{lemma-unibranch}{{0BQ0}{source}{}{sourcepointer-2}{}}
\newlabel{lemma-unibranch@cref}{{[originallemma][0][]0BQ0}{[1][0][]source}{}{}{}}
\makeatother
\crefname{originallemma}{original lemma}{original lemmas}
\Crefname{originallemma}{Original lemma}{Original lemmas}
\makeatletter
\newlabel{lemma-henselization-reduced}{{06DH}{source}{}{sourcepointer-3}{}}
\newlabel{lemma-henselization-reduced@cref}{{[originallemma][0][]06DH}{[1][0][]source}{}{}{}}
\makeatother
\crefname{originallemma}{original lemma}{original lemmas}
\Crefname{originallemma}{Original lemma}{Original lemmas}
\makeatletter
\newlabel{algebra-lemma-flat-going-down}{{00HS}{source}{}{sourcepointer-4}{}}
\newlabel{algebra-lemma-flat-going-down@cref}{{[originallemma][0][]00HS}{[1][0][]source}{}{}{}}
\makeatother
\crefname{originallemma}{original lemma}{original lemmas}
\Crefname{originallemma}{Original lemma}{Original lemmas}
\makeatletter
\newlabel{algebra-lemma-quotient-strict-henselization}{{05WS}{source}{}{sourcepointer-5}{}}
\newlabel{algebra-lemma-quotient-strict-henselization@cref}{{[originallemma][0][]05WS}{[1][0][]source}{}{}{}}
\makeatother
\crefname{originallemma}{original lemma}{original lemmas}
\Crefname{originallemma}{Original lemma}{Original lemmas}
\makeatletter
\newlabel{algebra-lemma-integral-going-up}{{00GU}{source}{}{sourcepointer-6}{}}
\newlabel{algebra-lemma-integral-going-up@cref}{{[originallemma][0][]00GU}{[1][0][]source}{}{}{}}
\makeatother
\crefname{originallemma}{original lemma}{original lemmas}
\Crefname{originallemma}{Original lemma}{Original lemmas}
\makeatletter
\newlabel{algebra-lemma-mop-up-strictly-henselian}{{06DD}{source}{}{sourcepointer-7}{}}
\newlabel{algebra-lemma-mop-up-strictly-henselian@cref}{{[originallemma][0][]06DD}{[1][0][]source}{}{}{}}
\makeatother
\crefname{originallemma}{original lemma}{original lemmas}
\Crefname{originallemma}{Original lemma}{Original lemmas}
\makeatletter
\newlabel{algebra-lemma-local-tensor-with-integral}{{092Y}{source}{}{sourcepointer-8}{}}
\newlabel{algebra-lemma-local-tensor-with-integral@cref}{{[originallemma][0][]092Y}{[1][0][]source}{}{}{}}
\makeatother
\crefname{originallemma}{original lemma}{original lemmas}
\Crefname{originallemma}{Original lemma}{Original lemmas}
\makeatletter
\newlabel{algebra-lemma-normal-goes-up}{{033C}{source}{}{sourcepointer-9}{}}
\newlabel{algebra-lemma-normal-goes-up@cref}{{[originallemma][0][]033C}{[1][0][]source}{}{}{}}
\makeatother
\title{A local ring is geometrically unibranch exactly when its strict henselization has a unique minimal prime}
\author{The Stacks Project Authors}
\date{stacks-project:more-algebra.tex:local-snapshot-20261012; GNU FDL 1.2 or later}
\begin{document}
\maketitle
\setcounter{subsection}{1}
\noindent\textit{Editorial note.} All delivered mathematical text is the author's own, copied from the pinned Stacks Project chapter file more-algebra.tex; no statement or proof step was rewritten, reordered or removed. The counted span is the author environment and its own proof environment. References inside the span are handed to the assembler as navigation stubs that carry the canonical tag of the cited tag only; whether each of them is load-bearing is recorded on the stub itself.\par\medskip
\noindent\textit{Allowed standard prerequisites (editorial).}\begin{itemize}
\item Flatness of the strict henselization makes the contraction of the unique minimal prime the unique minimal prime below, by going down. used at proof lines 31469-31472. Produces the unique minimal prime of the ring and the identification of the minimal prime of the strict henselization with its extension.
\item The quotient of the strict henselization by a prime is the strict henselization of the quotient and is reduced, so the minimal prime extends. invoked at proof lines 31472-31476. Identifies the minimal prime of the strict henselization as the extension of the prime below.
\item A finite local subalgebra of the integral closure obtained by adjoining finitely many elements keeps the relevant defect, and tensoring with the strict henselization turns it into a product of at least two local rings. used at proof lines 31482-31505. Produces the finite subalgebra B that leads to the contradiction with being a subring of a domain.
\item For a strictly henselian local ring, a finite algebra whose residue ring has two maximal ideals or a non purely inseparable residue field splits as a product. invoked at proof lines 31504-31507. Gives the splitting of the tensor product which contradicts its lying inside a domain.
\item Localizations of etale algebras preserve integrality, normality and the residue field conditions used in the converse. invoked at proof lines 31515-31520. Lets the converse pass from the integral closure of the ring to the tensor product with the localization of an etale algebra and conclude that the quotient is a domain.
\end{itemize}\medskip
\ifdefined\hypertarget\hypertarget{sourcepointer-1}{}\fi
\noindent\textit{Source reference.} Original lemma 0C25: sources/stacks-more-algebra/pinned-original-source.tex, source lines 31339--31339. The opening sentence notes the statement also follows from it, but the author gives the direct proof immediately; every step of that direct proof, including the contradiction with the splitting, is inside the counted span.\par\medskip
\ifdefined\hypertarget\hypertarget{sourcepointer-2}{}\fi
\noindent\textit{Source reference.} Original lemma 0BQ0: sources/stacks-more-algebra/pinned-original-source.tex, source lines 31249--31249. Cited at proof lines 31467-31469 only to say the direct proof is almost the same as the earlier one; the proof itself is written out in full inside the counted span.\par\medskip
\ifdefined\hypertarget\hypertarget{sourcepointer-3}{}\fi
\noindent\textit{Source reference.} Original lemma 06DH: sources/stacks-more-algebra/pinned-original-source.tex, source lines 11067--11067. Used at proof lines 31473-31475 to conclude that the quotient of the strict henselization by the prime below is reduced, which identifies the minimal prime above as the extension; the rest of the reduction is inside the counted span.\par\medskip
\ifdefined\hypertarget\hypertarget{sourcepointer-4}{}\fi
\noindent\textit{Source reference.} Original lemma 00HS: sources/stacks-chow/tags.txt, source lines 1--1. Applied at proof lines 31469-31472 to see that the contraction of the minimal prime above is the unique minimal prime below; the use of that uniqueness is inside the counted span.\par\medskip
\ifdefined\hypertarget\hypertarget{sourcepointer-5}{}\fi
\noindent\textit{Source reference.} Original lemma 05WS: sources/stacks-chow/tags.txt, source lines 1--1. Used at proof lines 31472-31474 to identify the quotient of the strict henselization by the prime with the strict henselization of the quotient; the conclusion drawn from that identification is inside the counted span.\par\medskip
\ifdefined\hypertarget\hypertarget{sourcepointer-6}{}\fi
\noindent\textit{Source reference.} Original lemma 00GU: sources/stacks-chow/tags.txt, source lines 1--1. Used at proof lines 31480-31482 to see that every maximal ideal of the integral closure lies over the maximal ideal of the ring; the construction of the finite subalgebra and the contradiction are inside the counted span.\par\medskip
\ifdefined\hypertarget\hypertarget{sourcepointer-7}{}\fi
\noindent\textit{Source reference.} Original lemma 06DD: sources/stacks-chow/tags.txt, source lines 1--1. Invoked at proof lines 31504-31507 to conclude that the tensor product with the strict henselization is a product of at least two local rings; producing the two maximal ideals of the residue tensor product is inside the counted span.\par\medskip
\ifdefined\hypertarget\hypertarget{sourcepointer-8}{}\fi
\noindent\textit{Source reference.} Original lemma 092Y: sources/stacks-chow/tags.txt, source lines 1--1. Used in the converse at proof lines 31517-31519 to see that the tensor product with the integral closure stays local; the surrounding argument that it is a normal domain and that the quotient is a domain is inside the counted span.\par\medskip
\ifdefined\hypertarget\hypertarget{sourcepointer-9}{}\fi
\noindent\textit{Source reference.} Original lemma 033C: sources/stacks-chow/tags.txt, source lines 1--1. Used at proof lines 31520-31523 to conclude that the tensor product is normal; the conclusion that it is a local normal domain and everything drawn from it is inside the counted span.\par\medskip
\begin{lemma}
\label{lemma-geometrically-unibranch}
\begin{reference}
\cite[Lemma 2.2]{Etale-coverings} and
\cite[Chapter IV Proposition 18.8.15]{EGA4}
\end{reference}
Let $A$ be a local ring. Let $A^{sh}$ be a strict henselization of $A$.
The following are equivalent
\begin{enumerate}
\item $A$ is geometrically unibranch, and
\item $A^{sh}$ has a unique minimal prime.
\end{enumerate}
\end{lemma}

\begin{proof}
This follows from Lemma \ref{lemma-geometric-branches}
but we will also give a direct proof; this direct
proof is almost exactly the same as the direct proof of
Lemma \ref{lemma-unibranch}.
Denote $\mathfrak m$ the maximal ideal of the ring $A$.
Denote $\kappa$, $\kappa^{sh}$ the residue field of $A$, $A^{sh}$.

\medskip\noindent
Assume (2). Let $\mathfrak p^{sh}$ be the unique minimal prime of
$A^{sh}$. The flatness of $A \to A^{sh}$ implies that
$\mathfrak p = A \cap \mathfrak p^{sh}$ is the unique minimal
prime of $A$ (by going down, see
Algebra, Lemma \ref{algebra-lemma-flat-going-down}).
Also, since $A^{sh}/\mathfrak pA^{sh} = (A/\mathfrak p)^{sh}$ (see
Algebra, Lemma \ref{algebra-lemma-quotient-strict-henselization})
is reduced by Lemma \ref{lemma-henselization-reduced}
we see that $\mathfrak p^{sh} = \mathfrak pA^{sh}$.
Let $A'$ be the integral closure of $A/\mathfrak p$ in its fraction
field. We have to show that $A'$ is local and that its residue
field is purely inseparable over $\kappa$.
Since $A \to A'$ is integral, every maximal ideal of $A'$ lies over
$\mathfrak m$ (by going up for integral ring maps, see
Algebra, Lemma \ref{algebra-lemma-integral-going-up}).
If $A'$ is not local, then we can find distinct maximal ideals
$\mathfrak m_1$, $\mathfrak m_2$. Choosing elements $f_1, f_2 \in A'$
with $f_i \in \mathfrak m_i, f_i \not \in \mathfrak m_{3 - i}$ we find
a finite subalgebra $B = A[f_1, f_2] \subset A'$ with distinct maximal
ideals $B \cap \mathfrak m_i$, $i = 1, 2$. If $A'$ is local with maximal
ideal $\mathfrak m'$, but $A/\mathfrak m \subset A'/\mathfrak m'$
is not purely inseparable, then we can find $f \in A'$ whose image in
$A'/\mathfrak m'$ generates a finite, not purely inseparable extension
of $A/\mathfrak m$ and we find a finite local subalgebra $B = A[f] \subset A'$
whose residue field is not a purely inseparable extension of $A/\mathfrak m$.
Note that the inclusions
$$
A/\mathfrak p \subset B \subset \kappa(\mathfrak p)
$$
give, on tensoring with the flat ring map $A \to A^{sh}$ the inclusions
$$
A^{sh}/\mathfrak p^{sh} \subset
B \otimes_A A^{sh} \subset
\kappa(\mathfrak p) \otimes_A A^{sh} \subset
\kappa(\mathfrak p^{sh})
$$
the last inclusion because
$\kappa(\mathfrak p) \otimes_A A^{sh} =
\kappa(\mathfrak p) \otimes_{A/\mathfrak p} A^{sh}/\mathfrak p^{sh}$
is a localization of the domain $A^{sh}/\mathfrak p^{sh}$.
Note that $B \otimes_A \kappa^{sh}$ has at least two maximal ideals
because $B/\mathfrak mB$ either has two maximal ideals or one whose
residue field is not purely inseparable over $\kappa$, and because
$\kappa^{sh}$ is separably algebraically closed. Hence, as
$A^{sh}$ is strictly henselian we see that
$B \otimes_A A^{sh}$ is a product of $\geq 2$ local rings, see
Algebra, Lemma \ref{algebra-lemma-mop-up-strictly-henselian}.
But we've just seen that $B \otimes_A A^{sh}$ is a subring of a domain
and we get a contradiction.

\medskip\noindent
Assume (1). Let $\mathfrak p \subset A$ be the unique minimal
prime and let $A'$ be the integral closure of $A/\mathfrak p$
in its fraction field. Let $A \to B$ be a local map of local rings which is a
localization of an \'etale $A$-algebra. In particular $\mathfrak m_B$
is the unique prime containing $\mathfrak m_AB$. Then $B' = A' \otimes_A B$
is integral over $B$ and the assumption that $A \to A'$ is local
with purely inseparable residue field extension implies that $B'$
is local (Algebra, Lemma \ref{algebra-lemma-local-tensor-with-integral}).
On the other hand, $A' \to B'$ is the localization
of an \'etale ring map, hence $B'$ is normal, see
Algebra, Lemma \ref{algebra-lemma-normal-goes-up}.
Thus $B'$ is a (local) normal domain. Finally, we have
$$
B/\mathfrak pB \subset B \otimes_A \kappa(\mathfrak p)
= B' \otimes_{A'} (\text{fraction field of }A') \subset
\text{fraction field of }B'
$$
Hence $B/\mathfrak pB$ is a domain, which implies that $B$ has a unique
minimal prime (since by flatness of $A \to B$ these all have to lie
over $\mathfrak p$). Since $A^{sh}$ is a filtered colimit of
the local rings $B$ it follows that $A^{sh}$ has a unique minimal prime.
Namely, if $fg = 0$ in $A^{sh}$ for some non-nilpotent elements
$f, g$, then we can find a $B$ as above containing both $f$ and $g$
which leads to a contradiction.
\end{proof}

\end{document}
